% !TeX root = ../../main.tex
\subsection{Endpoint benchmark and copy-number audit}

To ask which method ranks biological dependencies better, we used the
Brunello A375 negative-selection screen of \citet{sanson2018optimized}, bundled
with CB\textsuperscript{2}.  This is an ordinary endpoint screen rather than a
continuous-phenotype experiment, but it supplies an external target unavailable
in GSE70038: reference essential genes should rank ahead of reference
nonessential genes \citep{hart2017evaluation}.  The beta-binomial score combines guide
evidence by directional Stouffer aggregation; the MAGeCK score uses the
official MLE beta and Wald $p$-value.  Both are oriented so larger scores mean
stronger depletion.

\subsubsection{Copy-number-aware analysis}

Copy number is a gene-level property of A375, not a sample-level covariate:
for a given guide it is constant across the plasmid and three terminal
libraries.  Adding it directly to the four-row sample design would therefore
be non-identifiable with the endpoint effect.  Instead, we downloaded the
A375 gene-level profile (ACH-000219) from DepMap Public 19Q3
\citep{depmap19q3,meyers2017cnv} and used MAGeCK 0.5.9.5's official
\texttt{--cnv-norm} piecewise correction.  For an apples-to-apples comparison,
the same MAGeCK piecewise normalizer was applied to the beta-binomial
within-gene median effects.  It adjusts effect sizes after model fitting; it
does not refit either count likelihood.  The common CNV-complete evaluation set
contains 1,506 essential and 869 nonessential genes.

\begin{table}[ht]
\centering
\caption{CNV-aware Sanson A375 benchmark.  The final column is Spearman
correlation between effect and copy number among reference nonessential genes;
values nearer zero indicate less residual CNV-associated depletion.  F1 uses
nominal FDR 0.05 and a negative effect.  Higher discrimination metrics and
CNV correlation closest to zero are highlighted.}
\label{tab:sanson}
\small
\setlength{\tabcolsep}{4pt}
\begin{tabular}{lrrrr}
\toprule
Method & AUROC & Avg. precision & F1 & CNV correlation\\
\midrule
\methodswatch{barcsblue} BARCS &
  \textcolor{barcsblue}{\textbf{0.9598}} & 0.9786 &
  \textcolor{barcsblue}{\textbf{0.8531}} & $-0.185$\\
\methodswatch{barcsblue} BARCS + CNV &
  0.9533 & 0.9762 & \textcolor{barcsblue}{\textbf{0.8531}} & $-0.042$\\
\methodswatch{mageckorange} Official MAGeCK-MLE &
  \textcolor{mageckorange}{\textbf{0.9598}} &
  \textcolor{mageckorange}{\textbf{0.9795}} &
  0.8277 & $-0.207$\\
\methodswatch{mageckorange} Official MAGeCK-MLE + CNV &
  0.9501 & 0.9762 & 0.8277 &
  \textcolor{mageckorange}{$\mathbf{-0.006}$}\\
\bottomrule
\end{tabular}
\end{table}

CNV correction substantially removes the intended nuisance association:
$-0.185$ to $-0.042$ for beta-binomial effects and $-0.207$ to $-0.006$ for
MAGeCK.  It does not improve essential-gene discrimination in this screen.
AUROC decreases by 0.0066 and 0.0097, respectively, and the 0.05-FDR calls are
unchanged.  This is not a contradiction: bias removal and gold-standard
classification are different targets, and amplified genes can include both
dependencies and nondependencies.

After CNV correction, beta-binomial versus MAGeCK global ranking remains
indistinguishable: the AUROC difference is 0.00316 (paired
stratified-bootstrap 95\% CI $-0.00265$ to $0.00924$), and the average
precision difference is effectively zero (95\% CI $-0.00349$ to $0.00295$).
At nominal FDR 0.05, beta-binomial retains 4.18 percentage points more recall
(95\% CI 2.66 to 5.71) and an F1 advantage of 0.0253 (95\% CI 0.0148 to
0.0362).
